\section{検知条件の供給実態の定量調査}

\subsection{調査対象と方法}

調査対象は，我々が運用するデータ収集基盤が取得した次の2系統のスナップショットである．

\begin{description}
  \item[NVD] \mbox{} \\
    NVD JSON 2.0 データフィード（\texttt{nvdcve-2.0-*.json.gz}）由来，371,416件\cite{nvd-feed}
  \item[VulnCheck] \mbox{} \\
    VulnCheck NVD++ (\texttt{nist-nvd2}) API 由来，371,393件\cite{vulncheck-nvd2}
\end{description}

スナップショットは2026年7月29日時点（データセット中の \texttt{lastModified} の最大値）である．
各データソースは取得時のコミットハッシュで固定しており，付録\ref{app:repro}のとおり同一条件で再検証できる．
両者のCVE ID集合を比較すると，NVDのみに存在するものが23件，VulnCheckのみに存在するものが0件であった．
\textbf{VulnCheck NVD++ の収録IDはNVDのID空間を超えない}．
したがってVulnCheckは，5章で定義するエントリ欠落（C1）を原理的に改善しない．

以降の集計では \texttt{vulnStatus} が \texttt{Rejected} の17,712件を除外し，353,704件を対象とする．
取下げエントリは製品情報を持たないことが正常であり，含めると付与率を不当に低く見せるためである\footnote{実例として，Fortinet採番のCVEには2025年に「Rejected reason: Not used」として一括取下げされたものが866件含まれる．除外しない場合Fortinetの付与率は32.8\%と算出されるが，除外後は99.6\%である．}．

集計は，各エントリについて \texttt{configurations}（NVD）と \texttt{vcConfigurations}（VulnCheck独自）を走査し，\texttt{vulnerable: true} の \texttt{cpeMatch} の個数とバージョン指定の種別を数えることで行った．
測定スクリプトは付録\ref{app:repro}に従い再現できる．

\subsection{CPE適用条件の付与率の推移}

表\ref{tab:year}に公開年別の解析状況とCPE適用条件の付与率を示す．
「NVD付与率」は \texttt{configurations} に \texttt{vulnerable} な \texttt{cpeMatch} を1つ以上持つ割合，「VulnCheck付与率」は \texttt{configurations} または \texttt{vcConfigurations} のいずれかに持つ割合であり，「VC補完」はNVDが付与しておらずVulnCheckのみが付与している件数である．

\begin{table*}[htbp]
  \caption{公開年別のCPE適用条件の付与状況（取下げを除く）}
  \ecaption{CPE applicability coverage by publication year (rejected entries excluded).}
  \label{tab:year}
  \begin{center}
  \begin{tabular}{rrrrrrrr}
    \hline
    公開年 & CVE件数 & 解析完了\% & 解析待ち\% & Deferred\% & NVD付与率\% & VulnCheck付与率\% & VC補完件数（\%） \\
    \hline
    2020 & 18,322 & 100.0 & 0.0 & 0.0 & 100.0 & 100.0 & 0 (0.0) \\
    2021 & 20,149 & 100.0 & 0.0 & 0.0 & 100.0 & 100.0 & 0 (0.0) \\
    2022 & 25,074 & 100.0 & 0.0 & 0.0 & 100.0 & 100.0 & 0 (0.0) \\
    2023 & 28,817 & 100.0 & 0.0 & 0.0 & 100.0 & 100.0 & 0 (0.0) \\
    2024 & 39,957 & 80.5 & 0.0 & 19.5 & 80.5 & 88.3 & 3,083 (7.7) \\
    2025 & 48,162 & 62.6 & 0.0 & 37.3 & 62.6 & 76.1 & 6,492 (13.5) \\
    2026$^\dagger$ & 44,809 & 55.9 & 9.3 & 34.8 & 56.1 & 78.5 & 10,058 (22.4) \\
    \hline
  \end{tabular}
  \end{center}
  \footnotesize
  $\dagger$ 2026年は1月1日〜7月29日（210日間）の公開分．
  「解析完了」は \texttt{Analyzed} と \texttt{Modified} の合計，「解析待ち」は \texttt{Received}，\texttt{Awaiting Analysis}，\texttt{Undergoing Analysis} の合計．
\end{table*}

\paragraph{(1) 2023年公開分までは付与率100\%である．}
2016年から2023年までのいずれの年も，取下げを除く全エントリにCPE適用条件が付与されている．
「報告件数が激増する前のNVDは信頼できるのか」という問いに対し，少なくとも検知条件の\textbf{有無}という観点では，2023年公開分以前は完全に補われていると答えられる．
ただし記述内容が正しいかは別問題であり，先行研究\cite{dong2019,nguyen2013}が指摘するとおり一致率は別途評価が必要である（7章）．

\paragraph{(2) 2024年を境に急落する．}
付与率は2024年80.5\%，2025年62.6\%，2026年56.1\%と低下し，低下分はほぼそのまま \texttt{Deferred} の増加（19.5\%，37.3\%，34.8\%）で説明される．
付与率と解析完了率がほぼ一致することから，NVDにおけるCPE適用条件の付与はエンリッチメント完了とほぼ同義であることも確認できる．

\paragraph{(3) \texttt{Deferred} は待っても解消しない．}
その定義は「エンリッチメント予定に入っていない」であり，時間が経てば付与される保証はない．
2025年公開分は本稿執筆時点で公開から最短でも7か月，平均では約1年半が経過しているが，37.3\%が \texttt{Deferred} のままである．
すなわち2025年以降のNVDについては，\textbf{約4割弱は将来にわたって検知条件が供給されないと想定して設計すべき}である．

\paragraph{(4) VulnCheckは補完するが埋め切らない．}
独自エンリッチメントは付与率を2024年88.3\%，2025年76.1\%，2026年78.5\%へ引き上げ，補完件数は年を追って増加する（2026年は10,058件，22.4\%）．
商用エンリッチメントの寄与は無視できない大きさである．
一方で2025年以降は依然として2割強が両者ともに未付与である．

なお，VulnCheckは2024年11月の記事で，2024年公開分のCPE付与率を自社76.95\%に対しNVD 41.35\%と報告している\cite{vulncheck-cpe}．
同じ2024年公開分を本稿が2026年7月に観測すると，NVD 80.5\%・VulnCheck 88.3\%である．
集計定義が同一とは限らないため厳密な比較ではないが，\textbf{NVDは公開直後は大きく遅れるものの，1年半をかけて差の多くを埋める}（両者の差は35.6ポイントから7.8ポイントへ縮小）ことを示唆する．
一方で，そこから先は \texttt{Deferred} として残り続ける．
したがってVulnCheckの価値は，最終的な網羅性よりも公開直後の期間に集中すると解釈できる．

公開月別のコホート（付録\ref{app:cohort}）で見ると，公開から約2年経過した2024年9月公開分は解析完了88.5\%・\texttt{Deferred} 11.5\%，公開から約7か月の2025年12月公開分は解析完了50.4\%・\texttt{Deferred} 49.6\%，公開から2週間の2026年7月公開分は解析完了33.5\%・解析待ち30.4\%・\texttt{Deferred} 36.0\%であった．
時間経過とともに解析待ちは解消されるが，\texttt{Deferred} の割合は月によって25\%から50\%の幅で残る．

\paragraph{(5) 未付与エントリは長く滞留する．}
NVDがエンリッチメント対象と認識している \texttt{Awaiting Analysis} でも公開から中位数43日・90\%点124日が経過しており，\texttt{Deferred} の41,363件は中位数286日，うち32,215件が90日を超える（付録\ref{app:survey-detail}）．
\textbf{新規に公開されたCVEについては，公開後1〜4か月はNVDの検知条件が存在しない前提で設計する必要がある}．
これは，公開直後の脆弱性ほど攻撃に利用されやすいという性質と正面から衝突する．

\subsection{付与された条件の粒度とCNA別の偏り}

\textbf{CPE適用条件が存在することとバージョンを特定できることは同じではない}．
2024年以降に公開された132,928件のうち，バージョン範囲または特定バージョンで表現されているのはNVDで63.9\%，VulnCheckを加えて77.8\%にとどまる（付録\ref{app:survey-detail}）．
VulnCheckの \texttt{vcConfigurations} には両者が同一エントリ内に混在する例もある．
CVE-2024-6387（regreSSHion）では \texttt{cpe:2.3:a:openssh:openssh} に 8.5p1〜9.7p1 という適切な範囲がある一方，\texttt{cpe:2.3:a:gnu:grub2:-} のようにバージョン指定を持たない条件が並び，そのまま用いれば当該製品を持つ全ホストが検知対象となる．
\textbf{付与率の向上が検知精度の向上を意味するとは限らない}．

採番CNA別に見ると差は極端である（付録\ref{app:survey-detail}）．
Microsoft 99.7\%，Fortinet 99.6\%，Apple 99.4\%と付与率がほぼ維持される一方，Cisco 77.2\%（\texttt{Deferred} 15.0\%），\textbf{JPCERT/CC採番は39.1\%（\texttt{Deferred} 59.7\%）}でVulnCheckを加えても48.4\%にとどまる．
JPCERT/CCはJVNを通じて国内ベンダの脆弱性を採番する経路であり，\textbf{国内ベンダ製品の検知条件はNVDもVulnCheckもほぼ供給していない}．
NVDの \texttt{Deferred} は無作為ではなく領域に偏って適用されている．
またMicrosoftは付与率99.7\%に対しバージョン特定可能率が92.9\%と乖離する．
WindowsのCPEはバージョン部が \texttt{-} で表現され，パッチ適用状態は更新プログラム（KB）情報に依存するためである．

\subsection{未付与時に残る情報：CNA提供のaffected}

NVDがCPE適用条件を付与していない45,413件（2024年以降公開，取下げを除く）について代替情報の有無を調べた．
結果は，CNA提供の \texttt{affected} が存在するもの45,413件（100.0\%），VulnCheckの \texttt{vcConfigurations} が存在するもの19,607件（43.2\%），いずれも存在しないもの0件であった．

すなわち\textbf{製品情報そのものは常に存在し，欠けているのはCPEへの翻訳だけである}．
しかしその内容はそのまま照合に使えない．実例を示す．

\begin{description}
  \item[CVE-2025-20271] \mbox{} \\
    Cisco採番，2025-06-18公開，\texttt{Deferred}．\texttt{configurations} および \texttt{vcConfigurations} は空．\texttt{affected} は \texttt{\{"vendor": "Cisco", "product": "Cisco Meraki MX Firmware", "versions": [\{"version": "N/A", "status": "affected"\}]\}} のみで，バージョンが \texttt{N/A} である．一方Ciscoのアドバイザリ本体には修正済みバージョンが明記されている．
  \item[CVE-2023-39933] \mbox{} \\
    JPCERT/CC採番，2024-03-18公開，\texttt{Deferred}．\texttt{affected} は \texttt{\{"vendor": "A.K.I Software ", "product": "pmc.exe", "versions": [\{"version": "2.5.1.720 and earlier", "status": "affected"\}]\}}．バージョンが自然言語（\texttt{and earlier}）で記述され，ベンダ名には末尾空白が含まれる．
\end{description}

これらは，一次データソースへ遡り，自然言語のバージョン表記を機械可読な範囲式へ正規化する必要があることを示す．
6章の提案はこの要求に応える．
